\section{Stability near the symmetry that fixes the inert mode}
\label{app:houghton-stability}

We prove Theorem~\ref{thm:houghton-stability}. The round, clock torus,
constant fields and exterior-algebra representation are exactly those
of Appendix~\ref{app:houghton}; no change of geometry is made.
Write $\Pi(g)=(\bigwedge g)^{\otimes k}$, $L=4M+4$ and
$G_0=U(e_+^\perp)\cong U(L-1)$. All link permutations fix $e_+$,
as do the phase rotations on $(e_x-e_y)/\sqrt2$. Thus every literal
slot is a clock translation followed by a collective $G_0$ unitary.

Put $\sigma=\Delta_c\rho$, $\eta=\mathcal E_0\sigma$ and
$\delta=\frac12\|\sigma-\eta\|_1$. For clock-diagonal states,
$\mathcal E_0$ commutes with each literal conjugation: Haar twirling
is unchanged by the internal link on either side, and the same twirl
is used at every clock. Clock pinching also commutes with these
conjugations. The actual bath map is their weighted mixture.
Applying Holevo data processing to the ensemble of translated sources
therefore gives
\begin{equation}\label{eq:inert-twirl-gain}
 a(\eta)\leq a(\sigma)\leq a(\rho).
\end{equation}
This comparison does not assert that $\eta$ implements the full target
channel. We transfer only its order-three anchor charge.

\subsection{Entropy transferred from large irreducible sectors}

The one-particle three-cycle $\gamma$ fixes the orthogonal complement
of a complex two-plane in $e_+^\perp$ and has eigenvalues
$e^{2\pi i/3},e^{-2\pi i/3}$ on that plane. It is consequently an
element of the plane's $SU(2)$. In every representation, the projection
$Q_\gamma$ onto its nontrivial eigenvalues is bounded above by the
projection onto nontrivial $SU(2)$ summands.

For an irreducible $G_0$ representation $\lambda$, let $p_\lambda$ be
the flat-state mass of these nontrivial $SU(2)$ summands and
$q_\lambda$ the flat-state mass of $Q_\gamma$. Planes are conjugate in
$G_0$, so $p_\lambda$ is common to every plane and
$q_\lambda\leq p_\lambda$. Restrict to
$t=\lfloor(L-1)/2\rfloor=2M+1$ disjoint two-planes. Their commuting
$SU(2)$ groups decompose the flat irrep into flat tensor-product
irreducibles and multiplicities. A nontrivial spin factor has
dimension at least two. The entropy decomposition, including its
nonnegative mixing and multiplicity terms, gives
\begin{equation}\label{eq:inert-irrep-dimension}
 \log_2\dim\lambda\geq\sum_{j=1}^t\Pr(m_j>0)
                  =t p_\lambda\geq t q_\lambda.
\end{equation}
This is the same disjoint-plane argument as in
Appendix~\ref{app:houghton}, applied inside the smaller group.

The anchor calculation~\eqref{eq:actual-houghton-charge} gives
$\Tr\rho Q_\gamma\geq7/12$. It is unchanged by clock pinching.
Since $Q_\gamma$ is a projection, trace distance implies
$\Tr\eta Q_\gamma\geq7/12-\delta\geq1/2$ when $\delta\leq1/32$.
The clock-dependent three-cycles are conjugate in $G_0$, so this is
the average of the same charges $q_\lambda$ over the clock and
irrep weights of $\eta$. Irreps with $q_\lambda\geq1/4$ carry total
mass $w\geq1/3$, since the total charge is at most
$w+(1-w)/4$. By~\eqref{eq:inert-irrep-dimension}, each such irrep has
dimension at least $2^{M/2}$.

We use a dimension-free entropy transfer rather than continuity in
the ambient bath dimension.
\begin{lemma}[A large flat subdensity survives trace proximity]
\label{lem:flat-subdensity-transfer}
Let $\zeta,\xi$ be states, $\frac12\|\zeta-\xi\|_1\leq e$, and
$0\leq\xi_H\leq\xi$ with $\Tr\xi_H=w$ and
$\xi_H\leq2^{-B}I$. For every $s>0$,
\[
 S(\zeta)\geq s\bigl(w-e-2^{s-B}\bigr).
\]
\end{lemma}
\begin{proof}
Let $P$ be the spectral projection of $\zeta$ for eigenvalues at
least $2^{-s}$. Its rank is at most $2^s$. Hence
$\Tr\xi_H P\leq2^{s-B}$ and
$\Tr\zeta(I-P)\geq\Tr\xi(I-P)-e
\geq w-2^{s-B}-e$. Every eigenvalue of $\zeta$ on $I-P$ contributes
at least $s$ times its mass to the entropy. This also proves the
inequality when the right side is negative.
\end{proof}

Write $\sigma=\bigoplus_cp_c\sigma_c$ and
$\eta=\bigoplus_cp_c\eta_c$. Twirling preserves $p_c$.
Schur decomposition of a normalized invariant block gives
\[
 \eta_c=\bigoplus_\lambda
       \frac{I_{\dim\lambda}}{\dim\lambda}\otimes A_{c,\lambda},
 \qquad A_{c,\lambda}\succeq0,\quad
       \sum_\lambda\Tr A_{c,\lambda}=1.
\]
The high-charge subdensity is bounded above by $2^{-M/2}I$ even
when multiplicities are entangled or singular. Apply the lemma with
$B=M/2$, $s=M/4$ to each block, and average. The average high-charge
mass is $w\geq1/3$ and the average block trace distance is $\delta$.
Coherently copying the clock and applying Araki--Lieb, as in
Appendix~\ref{app:houghton}, gives
\begin{equation}\label{eq:inert-entropy-floor}
 S(\rho)\geq\sum_cp_c S(\sigma_c)
 \geq\frac M4\bigl(1/3-\delta-2^{-M/4}\bigr)
 \geq\frac M{16}\qquad(M\geq48,\ \delta\leq1/32).
\end{equation}
No physical clock measurement, full-rank hypothesis or eigenvalue
domination was used.

\subsection{Capped flux at its natural scale}

For an invariant internal state $\xi$, let $p_\xi$ be its common
nontrivial $SU(2)$-plane mass. The following estimate applies in
$G_0$ as well as in the full orbital group.
\begin{lemma}[Capped invariant flux]\label{lem:capped-invariant-flux}
Let $W$ have phases $\theta_j$ on disjoint one-particle modes in
$e_+^\perp$, and suppose these can be paired with distinct unused
modes in that space. If $|\theta_j|\leq\pi/4$ and
$\sum_j\theta_j^2\leq h\leq1$, then a collectively $G_0$-invariant
state satisfies
\[
 \Tr\xi\min\bigl(|\Pi(W)-I|^2,hI\bigr)
       \geq\frac{p_\xi}{2048}\sum_j\theta_j^2.
\]
The minimum denotes spectral functional calculus.
\end{lemma}
\begin{proof}
Let $C\in G_0$ swap the paired modes and set $H=WCW^*C^*$.
Its paired rotations belong to commuting $SU(2)$ groups. Conditional
on irreducible spins $m_j\geq0$, invariance makes their weights
independent and uniform:
\[
 X_j=(m_j-2r_j)\theta_j,\quad 0\leq r_j\leq m_j,
 \qquad v_j=\E X_j^2=m_j(m_j+2)\theta_j^2/3.
\]
They are symmetric, and their fourth moments are
$m_j(m_j+2)(3m_j^2+6m_j-4)\theta_j^4/15\leq3v_j^2$.
For $Z=\sum_jX_j$ and $v=\sum_jv_j$, independence gives
$\E Z^4\leq3v^2$.

If $0<v\leq1$, Paley--Zygmund gives
$\Pr(Z^2\geq v/2)\geq1/12$, while
$\Pr(|Z|>\pi)\leq3/\pi^4$. On the remaining event,
$4\sin^2(Z/2)\geq2v/\pi^2$. Thus
\[
 \E\min(4\sin^2(Z/2),h)
 \geq(1/12-3/\pi^4)(2/\pi^2)\min(v,h)
 >\min(v,h)/256.
\]
If $v\geq1$, each character $\varphi_j=\E e^{iX_j}$ satisfies
$|\varphi_j|\leq\exp[-\min(v_j,1)/64]$.
Indeed when $(m_j+1)|\theta_j|\leq1$, $|X_j|\leq1$ and the cosine
inequality gives $0\leq\varphi_j\leq e^{-2v_j/\pi^2}$.
Otherwise the spin character formula gives
\[
 |\varphi_j|=
 \left|\frac{\sin((m_j+1)\theta_j)}{(m_j+1)\theta_j}\right|
 \left|\frac{\theta_j}{\sin\theta_j}\right|
 \leq(101/120)(9/8)=909/960<e^{-1/64}.
\]
Here $|\sin z|/z\leq\sin1\leq101/120$ for $z\geq1$,
and $|\theta/\sin\theta|<9/8$ for $|\theta|\leq\pi/4$;
zero angles follow by continuity. Since
$\sum_j\min(v_j,1)\geq1$, the variable
$Y=4\sin^2(Z/2)\leq4$ has
$\E Y\geq2(1-e^{-1/64})$. The inequality
$\min(Y,h)\geq hY/4$ gives $\E\min(Y,h)>h/256$.
The case $v=0$ is immediate. In all cases the lower is
$\min(v,h)/256$. As
$v\geq\sum_{m_j>0}\theta_j^2$ and this sum is at most $h$,
averaging the spin sectors yields
\[
 \Tr\xi\min(|\Pi(H)-I|^2,hI)
                  \geq p_\xi\sum_j\theta_j^2/256.
\]

On the finite algebra generated by $\Pi(W),\Pi(C)$,
$\omega(A)=\Tr\xi A$ is tracial. For singular-value counting
$N(A;s)=\omega(1_{(s,\infty)}(|A|))$, the intersection/min--max
argument gives $N(A+B;s)\leq N(A;s/2)+N(B;s/2)$.
One can first remove zero-weight blocks if $\omega$ is not faithful.
Integration of this inequality gives
\[
 \omega\min(|A+B|^2,hI)
 \leq4\bigl[\omega\min(|A|^2,hI)+
             \omega\min(|B|^2,hI)\bigr].
\]
Use
\[
 \Pi(H)-I=
 [\, (\Pi(W)-I)\Pi(C)-\Pi(C)(\Pi(W)-I)\,]
                  \Pi(W)^*\Pi(C)^*.
\]
Tracial unitary invariance makes each summand's capped flux equal
to that of $\Pi(W)-I$, giving a factor eight. Combining the two
bounds proves the lemma, including arbitrary multiplicity states.
\end{proof}

In the actual geometry there are at most $M$ active modes. They lie
in $e_+^\perp$ and are orthogonal because their label pairs are
disjoint. Its dimension $4M+3$ leaves room for $M$ partners. With
$m_M=1+(M-1)(M-2)/2$, their constant-field angles satisfy
\[
 |\theta_j|\leq\pi/m_M\leq\pi/4,\qquad
 \sum_j\theta_j^2\leq M\pi^2/m_M^2\leq72/M^3=:h_M\leq1.
\]
The last estimate holds at $M=12$ and improves thereafter.
On $B_0=\{0,\ldots,\lfloor M/4\rfloor\}^2$, all $M$ regions
occur, with areas at most $M^2$, so
$\sum_j\theta_j^2\geq\pi^2/M^3$; $B_0$ occupies at least
$1/16$ of the clock torus.

Write $\eta_0=(I/M^2)\otimes\bar\eta$ for the uniform-clock state
with the same internal marginal as $\eta$. Its plane mass satisfies
$p_{\bar\eta}\geq\Tr\eta Q_\gamma\geq1/2$. Set
\[
 X=|\Pi(W)-I|^2,\quad F_M=\min(X,h_MI),\quad
 C_0=\pi^2/65536,
\]
where the clock-diagonal word $W=r_4$ uses the holonomy at each
clock. The capped lemma gives
$\Tr\eta_0F_M\geq C_0/M^3$.
The operator-valued clock law~\eqref{eq:operator-clock-law}
uses only invariance under the literal internal slots. It therefore
applies to $\eta$ with $G_0$ in place of $U(L)$, and
\eqref{eq:inert-twirl-gain} gives
\[
 \tfrac12\|\eta-\eta_0\|_1
           \leq M\sqrt{104(\ln2)a(\rho)}.
\]
Since $0\leq F_M\leq h_MI$ and its clock pinching is unchanged,
positivity and trace distance imply
\begin{equation}\label{eq:trace-stable-flux}
 \Tr\rho X\geq\Tr\rho F_M\geq
 M^{-3}\left[C_0-72\delta-
                 72M\sqrt{104(\ln2)a(\rho)}\right].
\end{equation}
Capping is important: the unbounded relative spectral contrast of
$\rho$ does not multiply the error by the ambient bath dimension.

\subsection{The cost bound and its limitation}

Put $\delta_*=C_0/288$ and
$A_0=\delta_*^2/(104\ln2)$. Assume $\delta\leq\delta_*$,
which is smaller than $1/32$. If $a>A_0/M^2$,
\eqref{eq:inert-entropy-floor} gives
\[
 S+na\geq M/16+nA_0/M^2
 \geq\frac3{2^{2/3}}(A_0/16^2)^{1/3}n^{1/3}
 >4.64\cdot10^{-6}n^{1/3}.
\]
Otherwise~\eqref{eq:trace-stable-flux} gives
$\Tr\rho X\geq C_0/(2M^3)$. The weaker consequence
$\ell\geq\Tr\rho X/(3\cdot104)$ of
\eqref{eq:actual-houghton-leakage}, together with
$\ell\leq2/(15n)$, implies
\[
 M^3\geq\frac{15C_0}{12\cdot104}n,\qquad
 S+na\geq\frac1{16}
       \left(\frac{15C_0}{12\cdot104}\right)^{1/3}n^{1/3}
 >7.61\cdot10^{-4}n^{1/3}.
\]
Both branches prove the stated $4\cdot10^{-6}$ constant.

The hypothesis is not density domination in disguise. Under the
canonical exterior-algebra splitting
$\Lambda(\C^L)\simeq\Lambda(\C e_+)\otimes\Lambda(e_+^\perp)$,
the first factor is inert for every literal slot. It can be in a pure
coherent state while the second carries an invariant charged source.
Such states can satisfy the theorem's round tests, have zero
eigenvalues, and fail full $U(L)$ invariance. For example, a flat
state on $\Lambda(e_+^\perp)^{\otimes k}$ has the same holonomy
characters as the flat full-orbital source; replacing its inert
factor by a pure coherent state changes no channel or leakage test.

More generally let $\tau$ have invariant clock-diagonal blocks and
let $\rho=P\tau P/\Tr(P\tau)$, where $\Tr(P\tau)=1-r$.
The purification proof of gentle filtering gives
$\frac12\|\rho-\tau\|_1\leq\sqrt r$. Pinching and twirling are
trace-distance contractions, so
$\delta_{\rm orb}(\rho)\leq2\sqrt r$. Whenever the original
round tests still hold, a vanishing filtering perturbation cannot
beat the exponent in this geometry. The missing all-source step is
to deduce a suitably uniform symmetry proximity from low actual
production as $M$ grows. No such deduction, or comparison with all
round geometries, is proved here.
